\section{Averaging as an Approximate Weighted LMO}
\label{app:matnormal-03}

The single-batch LMO in \Cref{sec:sample-loss-lmo} has two inputs, the
gradient in its objective and the per-sample second moment in its
outergradient constraint. To draw on recent batches while keeping the LMO
local to the current iterate, we evaluate every recent sample's gradient at
$W_t$ and average the two inputs across those samples, giving the ideal
moments $\overline G_t$ and $\overline\Sigma_t$. This appendix explains how
\Cref{alg:ours}
approximates them from the LoRA factor gradients already computed by
backpropagation. The load-bearing online fact is an exact bookkeeping identity:
the stored curvature EMA equals a current-closure fit plus a weighted
closure-drift term and the initialization tail. The rest of the appendix
explains what that current-closure fit estimates, and which approximation gaps
remain.

\Cref{sec:weighted-lmo-03} writes the weighted LMO and isolates the two
moments through which past steps enter the problem.
\Cref{sec:weighted-estimation-03}
first states the online identity, then separates the stale-gradient
approximation from the factor-gradient fit and connects the result to the
update~\eqref{eq:polora-update}.
\Cref{sec:scale-placement-03} explains why \Cref{alg:ours} normalizes the
accumulated vectors after averaging rather than normalizing each estimate
before it enters the average. The update is invariant to the scale of the
curvature factors, but the average that estimates them depends on that scale.

\subsection{The Weighted LMO}
\label{sec:weighted-lmo-03}

Let $s$ index optimizer steps, and let $i$ index samples in the
minibatch of size $n_s$ at step $s$. Write
$G_{s,i}:=\nabla\ell_{x_{s,i}}(W_s)$,
$g_{s,i}:=\vec(G_{s,i})$, and
$G_s:=n_s^{-1}\sum_iG_{s,i}$. Let $m_{t,s}$ and $w_{t,s}$ be
nonnegative weights on past steps, one for the objective and one for
the constraints.

In the ideal weighted LMO, every sample from a step $s\le t$ contributes
the gradient it would have at the current weights $W_t$, so drawing on past
steps changes only which samples enter the averages, not where their
gradients are evaluated. The two resulting moments are
\[
\overline G_t
  =
  \sum_{s\le t}\frac{m_{t,s}}{n_s}\sum_{i=1}^{n_s}
  \nabla\ell_{x_{s,i}}(W_t),
\qquad
\overline\Sigma_t
  =
  \sum_{s\le t}\frac{w_{t,s}}{n_s}\sum_{i=1}^{n_s}
  \vec(\nabla\ell_{x_{s,i}}(W_t))
  \vec(\nabla\ell_{x_{s,i}}(W_t))^\top .
\]
The first moment is the gradient in the linear objective. For the
constraint, the relevant gradients are rescaled as
$\sqrt{w_{t,s}/n_s}\,\nabla\ell_{x_{s,i}}(W_t)$, and their outer
products sum to $\overline\Sigma_t$. Applying the leverage argument of
\Cref{lem:gradprop} to these rescaled gradients gives a relaxation of
the outergradient constraint with the same form as~\eqref{eq:polorafull-hard},
with the scalar from the weighted average absorbed into the budget. The
relaxed problem is
\begin{equation}\label{eq:weighted-relaxed-03}
 \begin{array}{ll}
  \underset{\Delta W}{\mathrm{minimize}} &
  \ip{\overline G_t, \Delta W} \\[2pt]
  \mbox{subject to} &
  \displaystyle
  \max_{\widetilde{G}\in\mathcal{G}(\overline\Sigma_t)}
  \ip{\widetilde{G},\Delta W}\le\tau .
  \end{array}
\end{equation}
\Cref{alg:ours} stores neither directly, because it never recomputes old
samples' gradients at $W_t$ and the per-sample second moment is too large
to keep.

\paragraph{The exponential weights.}
If every past gradient could be recomputed at $W_t$, equal weights would be
the direct extension of the single-batch LMO to samples from earlier steps.
Since that average is unavailable, \Cref{alg:ours} uses two exponentially
decaying weight sequences, one for the objective and one for the constraint.
These weights use constant memory and define the weighted stale-gradient
discrepancies below. A current-batch-only estimate would discard the weighted
LMO target, while a uniform all-past observed-gradient estimate would keep old
stale gradients at full weight. The exponential average is the constant-memory
recency-weighted compromise.

\Cref{alg:ours} uses look-ahead momentum for the first moment,
\[
  M_{t+1}=\beta_1M_t+(1-\beta_1)G_t,
  \qquad
  \widehat M_t=\beta_1M_{t+1}+(1-\beta_1)G_t,
\]
and an EMA for the second moment. Unrolling these recursions
gives the induced step weights
\begin{equation}\label{eq:EMAweights-03}
  m_{t,s}=\beta_1^2(1-\beta_1)\beta_1^{t-1-s}
  \quad(s<t),
  \qquad
  m_{t,t}=1-\beta_1^2,
  \qquad
  w_{t,s}=(1-\beta_2)\beta_2^{t-s}.
\end{equation}
We write $\EMA_{s\le t}[\phi(G_s)]=\sum_{s\le t}w_{t,s}\,\phi(G_s)$ for an
average over past steps with weights $w_{t,s}$. Replacing each unavailable
current-weight gradient in $\overline G_t$ by the observed batch gradient
$G_s$ turns the first moment into the look-ahead momentum $\widehat M_t$,
with the objective weights $m_{t,s}$ in~\eqref{eq:EMAweights-03}. The
optimizer never forms the $d_{\mathrm{out}}\times d_{\mathrm{in}}$ matrix
$\widehat M_t$, so the next subsection expresses its needed projections and
the curvature EMA in terms of factor gradients.

\subsection{Estimating the Moments from Factor Gradients}
\label{sec:weighted-estimation-03}

The optimizer maintains momenta for $G_{A,s},G_{B,s}$ and diagonal
curvature vectors instead of the dense moments $\overline G_t$ and
$\overline\Sigma_t$.

\paragraph{The online recurrence before the model.}
The exact EMA identity does not require a Kronecker model. Let
$K^A_s,K^B_s$ be the symmetric inverse closures used when the factor gradients
at step $s$ are scored. \Cref{alg:ours} computes
\begin{equation}\label{eq:ours-klcoupling-03}
\begin{aligned}
  \hat q_s
    &= \frac1r\diag\!\left(G_{A,s}^\top K^A_sG_{A,s}\right),
  &
  \hat p_s
    &= \frac1r\diag\!\left(G_{B,s}K^B_sG_{B,s}^\top\right),
  \\
  q_{s+1}
    &= \beta_2q_s+(1-\beta_2)\hat q_s,
  &
  p_{s+1}
    &= \beta_2p_s+(1-\beta_2)\hat p_s .
\end{aligned}
\end{equation}
At time $t$, compare the stored EMA with the fit that scores the same past
factor gradients using the current closures:
\[
\begin{aligned}
  q^{\mathrm{cur}}_t
  &=
  \frac1r\diag\!\left(
    \sum_{s\le t}w_{t,s}G_{A,s}^\top K^A_tG_{A,s}
  \right),
  \\
  p^{\mathrm{cur}}_t
  &=
  \frac1r\diag\!\left(
    \sum_{s\le t}w_{t,s}G_{B,s}K^B_tG_{B,s}^\top
  \right).
\end{aligned}
\]
Unrolling~\eqref{eq:ours-klcoupling-03} and adding and subtracting these
current-closure sums gives
\begin{equation}\label{eq:online-current-closure-03}
\begin{aligned}
  q_{t+1}
  &=
  \beta_2^{t+1}q_0+q^{\mathrm{cur}}_t+d^A_t,\\
  d^A_t
  &=
  \frac1r\diag\!\left(
    \sum_{s\le t}w_{t,s}
    G_{A,s}^\top(K^A_s-K^A_t)G_{A,s}
  \right),\\
  p_{t+1}
  &=
  \beta_2^{t+1}p_0+p^{\mathrm{cur}}_t+d^B_t,\\
  d^B_t
  &=
  \frac1r\diag\!\left(
    \sum_{s\le t}w_{t,s}
    G_{B,s}(K^B_s-K^B_t)G_{B,s}^\top
  \right).
\end{aligned}
\end{equation}
This is the online estimator's approximation boundary. The remaining paragraphs
identify when the current-closure vectors have the diagonal shape predicted by
the dense curvature model.

\paragraph{From moments at $W_t$ to momentum and curvature variables.}
The ideal LMO in~\eqref{eq:weighted-relaxed-03} is written with old samples
re-evaluated at $W_t$. \Cref{alg:ours} replaces those ideal moments by stored
variables through the following modeling steps. Only the first step is a
stale-gradient approximation; the remaining steps concern the observed recent
gradient stream and the changing LoRA factor maps. The stale-gradient
approximation gives meaning to $\overline G_t$ and $\overline\Sigma_t$ as
current-weight LMO inputs; it does not by itself justify the online $p,q$
recurrence.
\begin{enumerate}[leftmargin=1.6em,itemsep=2pt,topsep=2pt]
  \item A gradient of an old sample at $W_t$ is replaced by the gradient
  observed when that sample was processed,
  $\nabla\ell_{x_{s,i}}(W_t)\leadsto\nabla\ell_{x_{s,i}}(W_s)$.
  This is the locality premise behind the recency weighting: over the mass of
  the weights $m_{t,s}$ and $w_{t,s}$, the samples that materially affect the
  average are recent enough that their gradients at $W_s$ are used as local
  surrogates for their gradients at $W_t$. The rest of the section works with
  the observed gradients $G_s$ and keeps this stale-gradient replacement as an
  approximation boundary, not an exact identity.
  \item The second moment within the batch is approximated by the outer
  product of the batch gradient,
  \[
    \frac1{n_s}\sum_i g_{s,i}g_{s,i}^\top
    \quad\leadsto\quad
    g_sg_s^\top,
    \qquad g_s=\vec(G_s).
  \]
  For $n_s>1$, this replaces the per-sample second moment by the
  second moment of the minibatch gradient, as in AdaGrad~\cite{adagrad},
  Adam~\cite{adam}, and Shampoo~\cite{shampoo}.
  \item When relating the dense observed moment to factor-gradient moments,
  replacing $B_s^\top G_s$ by $B_t^\top G_s$ and
  $G_sA_s^\top$ by $G_sA_t^\top$ is the factor-map approximation. The online
  statement below starts from the actually observed factor gradients and does
  not assume this replacement is exact.
  \item Assume the resulting local second moment has diagonal Kronecker factors
  $P^\star_t,Q^\star_t$,
  \begin{equation}\label{e-diag-kron-second-mom-model-03}
    \EMA_{s\le t}[g_sg_s^\top]
    =
    \kappa_t(Q^\star_t\otimes P^\star_t).
  \end{equation}
  We fix $\norm{\diag P^\star_t}_\infty=\norm{\diag Q^\star_t}_\infty=1$ and
  let the scalar $\kappa_t>0$ carry the magnitude removed by the final
  normalization.
\end{enumerate}
The stale-gradient and batch-gradient approximations reduce the constraint
moment to
\[
  \overline\Sigma_t\approx\EMA_{s\le t}[g_sg_s^\top],
\]
which the diagonal Kronecker model~\eqref{e-diag-kron-second-mom-model-03}
writes as $\kappa_t(Q^\star_t\otimes P^\star_t)$.

\paragraph{Factor momenta.}
LoRA backpropagation already computes the factor gradients
$G_{A,s}=B_s^\top G_s$ and $G_{B,s}=G_sA_s^\top$. Forming the dense
gradient $G_s$ would require a $d_{\mathrm{out}}\times d_{\mathrm{in}}$
matrix per layer. The optimizer therefore represents the first moment only
through the two projections needed by the factor update. Under the
linearized product update~\eqref{eq:DeltaW-lin}, the objective of
\eqref{eq:weighted-relaxed-03} decomposes as
\[
  \ip{\widehat M_t,\,\Delta B\,A+B\,\Delta A}
  =
  \ip{\widehat M_tA^\top,\Delta B}
  +
  \ip{B^\top\widehat M_t,\Delta A}.
\]
The direction step then needs only $\widehat M_tA_t^\top$ and
$B_t^\top\widehat M_t$. The factor momenta in \Cref{alg:ours} average
$G_sA_s^\top$ and $B_s^\top G_s$ with the objective weights $m_{t,s}$.
Reading these two EMAs as projections of the dense momentum at the current
factors uses the same factor-map approximation as above, now with the objective
weights $m_{t,s}$ rather than the curvature weights $w_{t,s}$.

\paragraph{Moments of the factor gradients.}
With $\EMA_{s\le t}[g_sg_s^\top]$ as the dense second moment, replacing
$A_s,B_s$ by $A_t,B_t$ links the second moments of the factor gradients to
the weighted second moment of $G_s$. For the $A$ factor,
$\vec(G_{A,s})=(I\otimes B_s^\top)\vec(G_s)$, so
\begin{align*}
 \Sigma_{A,t}
 &:= \EMA_{s\le t}\big[\vec(G_{A,s})\vec(G_{A,s})^\top\big] \\
 &= \EMA_{s\le t}\big[
      (I\otimes B_s^\top)g_sg_s^\top(I\otimes B_s)
    \big] \\
 &\approx
    (I\otimes B_t^\top)
    \EMA_{s\le t}\big[g_sg_s^\top\big]
    (I\otimes B_t) \\
 &\approx
    \kappa_t\,(I\otimes B_t^\top)
    (Q^\star_t\otimes P^\star_t)
    (I\otimes B_t) \\
 &= \kappa_t\,Q^\star_t\otimes(B_t^\top P^\star_tB_t).
\end{align*}
The first approximation replaces $B_s$ by $B_t$ inside the weighted
sum, the second applies the diagonal Kronecker model, and the final
equality is the mixed product rule. The analogous calculation for
$G_{B,s}$, written with $\vec(G_{B,s}^\top)$ so that the Kronecker
blocks are indexed by the output coordinate, gives
\begin{align}
  \Sigma_{A,t}
    &\approx \kappa_t\,Q^\star_t\otimes C^\star_{A,t},
  \nonumber \\
  \Sigma_{B,t}
    &:= \EMA_{s\le t}\big[
      \vec(G_{B,s}^\top)\vec(G_{B,s}^\top)^\top
    \big]
    \approx \kappa_t\,P^\star_t\otimes C^\star_{B,t},
  \label{eq:observable-moments-03}
\end{align}
where
$C^\star_{A,t}=B_t^\top P^\star_tB_t$ and
$C^\star_{B,t}=A_tQ^\star_tA_t^\top$.

\paragraph{Fitting the diagonal factors.}
We fit the diagonal factors to the factor-gradient moments in
\eqref{eq:observable-moments-03}. In this paragraph, all quantities are
at the current step, so we suppress the subscript $t$. Following
KL-Shampoo~\cite{klshampoo}, the fit uses the log-determinant
divergence~\cite{logdetdiv}
\[
  D(\Sigma,S)=\Tr(\Sigma S^{-1})-\log\det(\Sigma S^{-1})-d,
\]
defined for positive definite matrices of common dimension $d$.
Under the block model below, each factor-gradient moment separates a diagonal
coordinate factor from a small $r\times r$ closure. The fit freezes the
closure and estimates only the diagonal factor:
\begin{equation}\label{eq:factor-fit-03}
  \hat{q}(\widehat C_A)
  =
  \argmin_{q>0}
  D\bigl(\Sigma_A,\diag(q)\otimes\widehat C_A\bigr),
  \qquad
  \hat{p}(\widehat C_B)
  =
  \argmin_{p>0}
  D\bigl(\Sigma_B,\diag(p)\otimes\widehat C_B\bigr),
\end{equation}
where $\widehat C_A,\widehat C_B\succ0$. The blockwise first-order
conditions give the closed forms
\[
  \hat{q}(\widehat C_A)
  =
  \frac1r\diag\!\left(
    \EMA_{s\le t}[G_{A,s}^\top\widehat C_A^{-1}G_{A,s}]
  \right),
  \qquad
  \hat{p}(\widehat C_B)
  =
  \frac1r\diag\!\left(
    \EMA_{s\le t}[G_{B,s}\widehat C_B^{-1}G_{B,s}^\top]
  \right).
\]

\begin{lemma}[Scalar effect of the frozen closure]\label{lem:closure-invariance-03}
Assume the moments of the factor gradients satisfy
\[
  \Sigma_A=\kappa\,Q^\star\otimes C^\star_A,
  \qquad
  \Sigma_B=\kappa\,P^\star\otimes C^\star_B
\]
exactly, with
$C^\star_A=B^\top P^\star B$,
$C^\star_B=AQ^\star A^\top$, and
$\norm{\diag P^\star}_\infty=\norm{\diag Q^\star}_\infty=1$, with all
diagonal entries of $P^\star,Q^\star$ positive. Assume also that
$C^\star_A$ and $C^\star_B$ are positive definite.
Replacing the modeled closures by any positive definite closures
$\widehat C_A$ and $\widehat C_B$ changes the fitted vectors only by
side-specific scalars:
\[
  \hat{q}(\widehat C_A)
  =
  \kappa\frac{\Tr(C^\star_A\widehat C_A^{-1})}{r}\,\diag(Q^\star),
  \qquad
  \hat{p}(\widehat C_B)
  =
  \kappa\frac{\Tr(C^\star_B\widehat C_B^{-1})}{r}\,\diag(P^\star).
\]
Consequently $\hat{q}/\norm{\hat{q}}_\infty=\diag(Q^\star)$ and
$\hat{p}/\norm{\hat{p}}_\infty=\diag(P^\star)$.
\end{lemma}
\begin{proof}
Under the stated moment model, $\Sigma_A$ is block diagonal
with $j$th block $\kappa q^\star_j C^\star_A$, while
$\diag(q)\otimes\widehat C_A$ has $j$th block $q_j\widehat C_A$.
Traces and log-determinants add across blocks, so the objective for
$q$ is
\[
  \sum_j D(\kappa q^\star_j C^\star_A, q_j\widehat C_A).
\]
The $j$th term is minimized by its best scalar multiple,
\[
  \hat q_j
  =
  \kappa q^\star_j
  \frac{\Tr(C^\star_A\widehat C_A^{-1})}{r}.
\]
The scalar multiplier is independent of $j$, and the normalization removes it.
The proof for $\hat p$ is the same with the two factors exchanged.
\end{proof}

\paragraph{Closures used by the online estimator.}
The identity~\eqref{eq:online-current-closure-03} holds for any symmetric
closures $K^A_s,K^B_s$. The fixed-fit lemma applies to the current-closure
vectors $q^{\mathrm{cur}}_t,p^{\mathrm{cur}}_t$, not directly to the stored EMA.
This paragraph instantiates the generic closures with the quantities used by
\Cref{alg:ours}. At step $t$, these closures are formed before the factor update
and then reused in the curvature update at the end of the step. Given raw EMA
vectors $p_t,q_t$, the implementation first forms the normalized diagonal
factors in \Cref{alg:ours},
\[
  P_t=\diag\!\bigl(p_t/\norm{p_t}_\infty\bigr),
  \qquad
  Q_t=\diag\!\bigl(q_t/\norm{q_t}_\infty\bigr).
\]
The raw vectors $p_t,q_t$ are initialized with positive entries, so these
normalizations are defined before the first curvature update. The displayed
algorithm writes the estimator with $C_A^{-1}$ and $C_B^{-1}$. The shipped
implementation applies the relative damping convention at two places. First it
replaces the normalized diagonal metrics by
\[
  \widetilde P_t
  =
  \diag\!\bigl(p_t/\norm{p_t}_\infty+\delta\mathbf 1\bigr),
  \qquad
  \widetilde Q_t
  =
  \diag\!\bigl(q_t/\norm{q_t}_\infty+\delta\mathbf 1\bigr).
\]
It then forms the small closures from these effective metrics and applies
relative damping again to the small inverse closures. In the exact inverse
idealization, this gives the representatives
\[
\begin{aligned}
  \widetilde C_{A,t}&=B_t^\top \widetilde P_tB_t,
  &
  \widetilde C_{B,t}&=A_t\widetilde Q_tA_t^\top,\\
  K^A_t
  &\propto(\widetilde C_{A,t}
    +\delta\lambda_{\max}(\widetilde C_{A,t})I)^{-1},
  &
  K^B_t
  &\propto(\widetilde C_{B,t}
    +\delta\lambda_{\max}(\widetilde C_{B,t})I)^{-1}.
\end{aligned}
\]
Here $\delta>0$ is the relative damping parameter. Iterative inverse or inverse
square-root implementations approximate this representative; that numerical
approximation is outside the estimator identity and can be absorbed into the
choice of $K^A_t,K^B_t$. The idealized displayed algorithm instead uses
$C_{A,t}=B_t^\top P_tB_t$, $C_{B,t}=A_tQ_tA_t^\top$ and
$K^A_t=C_{A,t}^{-1}$, $K^B_t=C_{B,t}^{-1}$. The proportionality records the
scale convention used by the damped inverse. In
\eqref{eq:online-current-closure-03}, $K^A_t,K^B_t$ mean whichever symmetric
representatives are used to score the curvature observation; any side-specific
scale change across time is part of the closure-drift term.

\paragraph{Stationary block-moment check.}\label{par:stationary-check-03}
Consider the tail-free, moment-level version of the online recurrence, where
each single-batch quadratic observation is replaced by the corresponding
factor-gradient moment at that step. Suppose that, for every $s\le t$ in the
full weighted sum, these moments share the same normalized diagonal block
factors:
\[
  \Sigma_{A,s}=\kappa^A_s\,Q^\star\otimes C^\star_{A,s},
  \qquad
  \Sigma_{B,s}=\kappa^B_s\,P^\star\otimes C^\star_{B,s},
\]
where $P^\star,Q^\star$ have positive diagonal entries and are normalized.
The scalars $\kappa^A_s,\kappa^B_s$ and the closure blocks
$C^\star_{A,s},C^\star_{B,s}$ may change with $s$. Let the recurrence score
step $s$ with any positive definite inverse closures $K^A_s,K^B_s$. For the
$A$ side, the $j$th block of $\Sigma_{A,s}$ is
$\kappa^A_s q^\star_j C^\star_{A,s}$, so scoring that block with $K^A_s$ gives
\[
  \hat q_{s,j}
  =
  \frac1r\Tr\!\left(\kappa^A_s q^\star_j C^\star_{A,s}K^A_s\right)
  =
  \alpha^A_s q^\star_j,
  \qquad
  \alpha^A_s:=\frac{\kappa^A_s}{r}\Tr(C^\star_{A,s}K^A_s)>0 .
\]
Thus the $A$-side tail-free moment-level EMA is
$(\sum_{s\le t}w_{t,s}\alpha^A_s)\diag(Q^\star)$. The $B$ side is identical,
with a scalar
$\alpha^B_s=r^{-1}\kappa^B_s\Tr(C^\star_{B,s}K^B_s)>0$. The final
normalization removes the positive scalars. This check uses a shared diagonal
block shape in the moment model, not stationarity of the already-computed
running estimates.

When the diagonal model factors are not shared across $s$, the same
moment-level calculation returns a scalar-weighted average of the recent
$\diag(Q^\star_s)$ and $\diag(P^\star_s)$. The actual stored recurrence also
contains the initialization tail and the sample-level discrepancy between
quadratic observations and their moment model. Thus the shared-shape condition
is on the block-moment model factors in the full weighted sum, not on the
already-computed running estimates.

\paragraph{Relation to KL-Shampoo.}
KL-Shampoo~\cite{klshampoo} also uses the log-determinant divergence to
fit Kronecker factors. The fit here differs in three ways.
\begin{itemize}[leftmargin=1.5em,itemsep=1pt,topsep=2pt]
  \item \emph{Finite EMA.} KL-Shampoo is framed around fitting moments of a
  gradient distribution. Here $\EMA_{s\le t}$ is a finite weighted sum over
  realized gradients, and the online comparison is the finite identity above.
  \item \emph{The fitted moments are factor-gradient moments.} KL-Shampoo
  fits Kronecker factors to the second moment of the weight gradient $G$.
  Here $G$ is never formed. Each diagonal factor is fit from the second
  moment of its corresponding LoRA factor gradient, and the two fits are
  coupled through $C_A$ and $C_B$.
  \item \emph{Fixed-moment closure invariance.}
  \Cref{lem:closure-invariance-03} and the stationary block-moment check are
  algebraic statements for block moments.
  Conditional on those block models, the closure choice only introduces
  side-specific scalars, which the normalization removes. The online EMA differs
  from the current fixed-closure fit by the weighted closure-drift terms
  displayed above.
\end{itemize}

\paragraph{Connection to the update of \Cref{alg:ours}.}
This paragraph maps the stored quantities to the implemented update. The
interpretation as an approximate weighted LMO uses the two bridges separately:
stale-gradient locality maps the ideal current-weight moments to observed
moments, and the closure-drift identity maps the current-closure fit to the
online $p,q$ recurrence.
With the diagonal Kronecker model used for the second moment and the
look-ahead momentum approximation used for the gradient input, the spectral LMO
has the preconditioned form in~\eqref{eq:polorafull}
(\Cref{lem:pseudogradient-spectral}). The direction and magnitude steps of
\Cref{sec:sample-loss-lmo}, using the factor momenta for the two dense momentum
projections, are the update~\eqref{eq:polora-update}. \Cref{alg:ours} updates
the curvature factors after the direction step, so step $t$ uses the estimate
from the previous step. This avoids forming and inverting $C_A$ and $C_B$ twice
per optimizer step.

\paragraph{What this establishes.}
The derivation gives a concrete target for the state maintained by
\Cref{alg:ours}. Under the weighted stale-gradient approximation, the
batch-gradient second-moment approximation, the factor-map approximation, and
the diagonal Kronecker model for the observed dense moment, the fixed-closure
fit has the normalized diagonal shape $Q^\star_t$ and $P^\star_t$.
The stored EMA is not identified with that fit by definition: its exact
difference from the current-closure fit is the weighted closure-drift vector
displayed above, plus the initialization tail. Thus \Cref{alg:ours} is an
online approximation to the weighted LMO in the following specific sense. The
factor momenta approximate the two dense momentum projections. The fixed-closure
fit has the model diagonal shape only under the exact block model, and the
stationary block-moment check shows one exact regime where changing closures
does not alter that normalized shape. Outside these exact statements, the
stale-gradient replacement, batch-gradient replacement, factor-map replacement,
diagonal Kronecker modeling, model-factor variation, and sample-level
finite-batch effects are unquantified approximation gaps. The only exact online
gap displayed here is the closure-drift identity together with the initialization
tail.

\subsection{Normalization Placement}
\label{sec:scale-placement-03}

\methodname{} stores raw accumulated vectors $q_{t+1},p_{t+1}$ and normalizes
them only when forming
\[
  Q_{t+1}=\diag(q_{t+1}/\norm{q_{t+1}}_\infty),
  \qquad
  P_{t+1}=\diag(p_{t+1}/\norm{p_{t+1}}_\infty).
\]
It does not normalize each fitted observation $\hat q_t,\hat p_t$ before
adding it to the EMA.

\paragraph{Raw scale does not change the implemented closures.}
For any $a,b>0$,
\[
  \diag(aq/\norm{aq}_\infty)=\diag(q/\norm{q}_\infty),
  \qquad
  \diag(bp/\norm{bp}_\infty)=\diag(p/\norm{p}_\infty).
\]
The rank-$r$ closures $C_A,C_B$ and the relative-damped inverse closures
$K^A,K^B$ are therefore unchanged by
positive rescaling of the raw EMA vectors. This is the implementation-level
scale invariance. The magnitude $\kappa_t$ in the diagonal Kronecker
model~\eqref{e-diag-kron-second-mom-model-03} is removed before the
direction step sees the curvature factors.

\paragraph{Normalizing each observation would change the estimator.}
The recurrence in~\eqref{eq:ours-klcoupling-03} estimates the weighted
arithmetic mean of the fitted observations,
\[
  q^\circ_{t+1}=\sum_{s\le t}w_{t,s}\hat q_s,
  \qquad
  p^\circ_{t+1}=\sum_{s\le t}w_{t,s}\hat p_s,
\]
with the raw vectors differing by the initialization terms in the unrolled
recurrence. Normalizing before the EMA would replace this by
\[
  \sum_{s\le t}w_{t,s}\frac{\hat q_s}{\norm{\hat q_s}_\infty},
  \qquad
  \sum_{s\le t}w_{t,s}\frac{\hat p_s}{\norm{\hat p_s}_\infty},
\]
which gives a small-magnitude and a large-magnitude fitted observation the
same maximum coordinate before averaging. The implemented order averages the
raw observations first and applies the projective normalization only when the
next effective metric is formed. The same average-then-normalize order appears
in normalized SGD~\cite{cutkosky_nigt}, where momentum precedes
normalization, and in Muon~\cite{denoise_ortho}, where momentum precedes
orthogonalization.
